\documentclass[10pt,a4paper]{amsart}
\usepackage[margin=19mm]{geometry}
\usepackage{amsmath,amssymb,mathtools}
\usepackage[scr=rsfs,cal=euler]{mathalfa}
\usepackage{xfrac}
\usepackage[unicode,hidelinks,bookmarks=false]{hyperref}
\newcommand{\ie}{i.e.\ }
\newcommand{\cf}{cf.\ }
\newcommand{\resp}{respectively\ }
\newcommand{\Subsection}{\subsection}
\providecommand{\nobreakdash}{\nobreak\hbox{-}}
\setlength{\emergencystretch}{2em}
\multlinegap0pt
\usepackage{amssymb}
\usepackage{mathtools}
\usepackage{xcolor}
\newtheorem{lemma}{Lemma}
\newtheorem{proposition}{Proposition}
\newcommand{\Aut}{\operatorname{Aut}}
\newcommand{\GN}{\Gamma_{\mathrm N}}
\newcommand{\VZG}{V(\mathbb ZG)}
\newcommand{\eps}{\varepsilon}
\title{The $N$-Prime Graph Question\\ is equivalent to the Prime Graph Question}
\author{Brecht Verbeken}
\date{Source: arXiv:2607.29105v1 (CC BY 4.0)}
\begin{document}
\maketitle

\noindent\emph{Source and licence.} The $N$-Prime Graph Question\\ is equivalent to the Prime Graph Question. Version arXiv:2607.29105v1 (CC BY 4.0), distributed by arXiv under the Creative Commons Attribution 4.0 International licence, http://creativecommons.org/licenses/by/4.0/, as recorded in the arXiv licence metadata for this exact version and shown on the abstract page https://arxiv.org/abs/2607.29105v1. Full licence notice: \texttt{licenses/arxiv-2607.29105-CC-BY-4.0.txt}.\par\smallskip

\noindent\textit{Editorial scope.} Counted unit: the article's proposition prop:transfer (source0.tex lines 268--276). The article's complete original proof of it is delivered with it (source0.tex lines 278--327, one \texttt{proof} environment of 50 lines). No mathematics is written, repaired or re-derived: the statement and the proof are the article's own lines, in the article's own notation, and no retained line is substituted or re-flowed.  Retained uncounted as the source-local context the proof needs: lines 215--219; lines 234--245.  Nothing was omitted from any retained span, so the package carries no omission record.  No retained line contains a citation, so the wrapper carries no bibliography and no bibliography entry.  No retained line contains a figure environment, an image-inclusion command or drawing code, so no image is delivered and the package declares no asset dependency.  Editorial formatting only, in addition: an AMS \texttt{amsart} preamble that restates the article's own declarations and macros used by the retained lines, namely the environments \texttt{lemma}, \texttt{proposition} on separate counters, together with the standard packages those lines need.  The retained environments print in this document as Lemma 1, Proposition 1 in order of appearance, whereas the article prints the same items with the numbering of its own full document.  The complete native source of the article as distributed in the arXiv e-print for this exact version is bundled unchanged as \texttt{sources/206460/source0.tex}.
\begin{equation}\label{eq:support}
 \eps_C(u)=0\quad(C\notin\mathcal C_p(G)),
 \qquad
 \sum_{C\in\mathcal C_p(G)}\eps_C(u)=1.
\end{equation}

\begin{lemma}\label{lem:power}
Let $u\in\VZG$ have order $p$, where $p$ is prime, and let $r$ be coprime
to $p$. Then
\[
 \eps_C(u^r)
 =
 \eps_{C^{r^{-1}}}(u)
 \qquad
 (C\in\mathcal C_p(G)),
\]
where $r^{-1}$ is taken modulo $p$.
\end{lemma}

\begin{proposition}[Transfer of nontrivial arcs]\label{prop:transfer}
Let $p$ and $q$ be distinct primes. Suppose that $u,v\in\VZG$ satisfy
\[
 |u|=p,\qquad |v|=q,\qquad v^{-1}uv=u^r,
 \qquad \operatorname{ord}_p(r)=q.
\]
Then the map $C\mapsto C^r$ fixes a conjugacy class of elements of order
$p$ in $G$. Consequently, $q\to p$ is an arc of $\GN(G)$.
\end{proposition}

\begin{proof}
Let
\[
 \rho\colon\mathcal C_p(G)\longrightarrow\mathcal C_p(G),
 \qquad \rho(C)=C^r.
\]
Since $\gcd(r,p)=1$, the map $\rho$ is a permutation. Conjugation
invariance of partial augmentations and Lemma~\ref{lem:power} give
\[
 \eps_C(u)
 =
 \eps_C(v^{-1}uv)
 =
 \eps_C(u^r)
 =
 \eps_{\rho^{-1}(C)}(u)
 \qquad(C\in\mathcal C_p(G)).
\]
Thus the partial augmentations of $u$ are constant on the $\rho$-orbits.
Moreover, $r^q\equiv1\pmod p$ implies $\rho^q=\operatorname{id}$. Hence
every $\rho$-orbit has size dividing the prime $q$, and therefore has size
$1$ or $q$. By \eqref{eq:support},
\[
 1
 =
 \sum_{C\in\mathcal C_p(G)}\eps_C(u)
 \equiv
 \sum_{\substack{C\in\mathcal C_p(G)\\ C^r=C}}
 \eps_C(u)
 \pmod q.
\]
In particular, the power map has a fixed class.

Let $a^G$ be such a class. Since $a^r$ is conjugate to $a$, choose $x\in G$
with $x^{-1}ax=a^r$. Since $a^r$ generates $\langle a\rangle$, we have
$x\in N_G(\langle a\rangle)$. Under the standard conjugation homomorphism
\[
 N_G(\langle a\rangle)\longrightarrow
 \Aut(\langle a\rangle),
 \qquad y\longmapsto(t\longmapsto yty^{-1}),
\]
the element $x^{-1}$ induces $a\mapsto a^r$. Its image therefore has order
$\operatorname{ord}_p(r)=q$. Thus
\[
 q\mid\bigl|N_G(\langle a\rangle)/C_G(a)\bigr|,
\]
and in particular $q\mid |N_G(\langle a\rangle)|$. Cauchy's theorem now
gives an element of order $q$ in this normalizer. Thus $q\to p$ is an arc
of $\GN(G)$.
\end{proof}

\end{document}
